% ============================================================
% PHẦN 7 — Slide 61-67 (Quan sát bổ sung, Tổng kết, Q&A)
% ============================================================

% --- Slide 61 ---
\begin{frame}{Quan sát 1: \texttt{final\_prune} không kịp chạy ở 7k iteration}
\begin{itemize}
    \item Hàm \texttt{final\_prune\_fastgs} chỉ được phép chạy trong cửa sổ bảo vệ:
    \[
        15000 < \text{iteration} < 30000
    \]
    \item Một lần chạy chỉ với \textbf{7000 iteration} \textbf{không bao giờ} chạm tới cửa sổ này $\Rightarrow$ mô hình giao ra \textbf{chưa từng được prune cuối cùng}.
    \item Đây là lời giải thích khả dĩ nhất cho:
    \begin{itemize}
        \item Kích thước file lớn bất thường: \textbf{248 B/Gaussian}, hoàn toàn chưa nén.
        \item LPIPS "chững" quanh \textbf{0.31} thay vì tiếp tục cải thiện.
    \end{itemize}
    \item \alert{Hệ quả quan trọng:} mọi so sánh trực tiếp giữa số liệu 7k-iteration này với kết quả công bố gốc của FastGS/3DGS (chạy đủ 30k iteration) đều \textbf{không hợp lệ} — khác ngân sách huấn luyện, khác điểm dừng pipeline.
\end{itemize}
\end{frame}

% --- Slide 62 ---
\begin{frame}{Quan sát 2: Nút thắt là RAM hệ thống, không phải VRAM}
\begin{itemize}
    \item Đo trên Google Colab T4 trong quá trình huấn luyện:
\end{itemize}
\footnotesize
\begin{table}[]
\centering
\begin{tabular}{lccc}
\toprule
Tài nguyên & Đỉnh sử dụng & Khả dụng & \% sử dụng \\
\midrule
VRAM (GPU) & 1.16 GB & 14.56 GB & \textbf{8\%} \\
RAM hệ thống & 9.59 GB & 12.7 GB & \textbf{91\%} (giới hạn mềm 10.5 GB) \\
\bottomrule
\end{tabular}
\end{table}
\normalsize
\begin{itemize}
    \item RAM hệ thống áp sát giới hạn mềm của pipeline, trong khi VRAM còn dư rất nhiều.
    \item Khuyến nghị cũ "giảm độ phân giải để tiết kiệm VRAM" \textbf{không cần thiết} ở ngân sách này.
    \item \texttt{resolution=1} (độ phân giải gốc) hoàn toàn khả thi và còn loại bỏ luôn độ lệch train/eval resolution.
\end{itemize}
\end{frame}

% --- Slide 63 ---
\begin{frame}{Quan sát 3: Live-score vs reported-score — khác protocol, không phải suy giảm}
\begin{itemize}
    \item Score đo \textbf{trong lúc train} (live) cao hơn hẳn score \textbf{báo cáo cuối} (reported) — nhưng đây không phải mô hình bị tệ đi.
\end{itemize}
\footnotesize
\begin{table}[]
\centering
\begin{tabular}{lccc}
\toprule
Metric & Live (trong train) & Reported (cuối) & Chênh lệch \\
\midrule
PSNR (dB) & 26.23 & 24.46 & $-1.77$ \\
SSIM & 0.8675 & 0.8142 & $-0.0533$ (ít nhất) \\
LPIPS & 0.114 & 0.311 & $+0.197$ \\
\bottomrule
\end{tabular}
\end{table}
\normalsize
\begin{itemize}
    \item Nguyên nhân: (1) LPIPS backbone khác nhau — AlexNet lúc train vs VGG lúc báo cáo; (2) độ phân giải khác nhau — nửa độ phân giải lúc train vs độ phân giải gốc lúc báo cáo.
    \item \alert{Bài học:} không so sánh chéo hai protocol đo khác nhau khi đánh giá tiến bộ giữa các lần chạy.
\end{itemize}
\end{frame}

% --- Slide 64 ---
\begin{frame}{Tổng kết pipeline render 3D Gaussian Splatting}
\centering
\footnotesize
\begin{tikzpicture}[
    node distance=8mm and 6mm,
    box/.style={draw, rounded corners, align=center, minimum height=0.9cm, minimum width=2.6cm, fill=blue!8}
]
\node[box] (gauss) {3D Gaussian\\$(\mu,\Sigma,\alpha,SH)$};
\node[box, right=of gauss] (proj) {Chiếu phối cảnh\\(EWA)};
\node[box, right=of proj] (raster) {Rasterizer khả vi\\theo tile};

\node[box, below=of raster] (blend) {Alpha\\blending};
\node[box, left=of blend] (loss) {Loss\\($L_1$+D-SSIM)};
\node[box, left=of loss] (back) {Backward\\(gradient)};

\node[box, below=of back, fill=orange!12] (adc) {Adaptive\\Density Control};

\draw[-Latex] (gauss) -- (proj);
\draw[-Latex] (proj) -- (raster);
\draw[-Latex] (raster) -- (blend);
\draw[-Latex] (blend) -- (loss);
\draw[-Latex] (loss) -- (back);
\draw[-Latex] (back) -- (adc);
\draw[-Latex] (adc) -- (gauss);
\end{tikzpicture}
\vspace{2mm}

\normalsize
Một chu trình khép kín: render $\rightarrow$ tính loss $\rightarrow$ lan truyền ngược $\rightarrow$ điều chỉnh mật độ Gaussian $\rightarrow$ lặp lại.
\end{frame}

% --- Slide 65 ---
\begin{frame}{Tổng kết cải tiến FastGS \& giới hạn còn lại}
\textbf{Ba đòn bẩy chính và tác động:}
\begin{itemize}
    \item \textbf{Multi-view score} + \textbf{Compact box} $\Rightarrow$ giảm số Gaussian dư thừa mà không giảm chất lượng đáng kể.
    \item \textbf{Sparse Adam} $\Rightarrow$ giảm chi phí tính toán mỗi iteration (chỉ cập nhật tham số liên quan).
    \item \textbf{Mới:} tích hợp thêm 5 cơ chế từ \textbf{Faster-GS} (Hahlbohm et al., CVPR 2026) — fused Adam, 3D anti-aliasing filter, Morton reordering, random-init fallback (luôn bật), MCMC densification (có sẵn, giữ tắt) — xem phần Faster-GS mới trong bộ slide này để biết chi tiết công thức.
\end{itemize}
\vspace{2mm}
\textbf{Giới hạn / hướng phát triển:}
\begin{enumerate}
    \item Dữ liệu Gaussian đầu ra chưa được nén (248 B/Gaussian, thô).
    \item Ngân sách 7k iteration của fastgs-lite bỏ lỡ bước \texttt{final\_prune} (xem Slide 61).
    \item Chưa có ablation study cô lập riêng từng đòn bẩy để đo đóng góp thực tế của mỗi thành phần.
    \item 4 nhánh mở rộng của FastGS gốc — dynamic scenes, sparse view, surface reconstruction, SLAM — chưa được tích hợp vào fork này.
    \item Các cơ chế Faster-GS mới tích hợp mới chỉ được review tĩnh (đọc code), \textbf{CHƯA} build/benchmark thật trên GPU (không có CUDA toolchain khi phát triển).
\end{enumerate}
\end{frame}

% --- Slide 66 ---
\begin{frame}{Bản đồ tài liệu tham khảo}
\footnotesize
\begin{table}[]
\centering
\begin{tabular}{p{4.2cm}p{6.5cm}}
\toprule
\textbf{Tài liệu} & \textbf{Nội dung} \\
\midrule
\texttt{README.md} & Điểm vào tiếng Anh, bảng kết quả đầy đủ \\
\texttt{DOCS/fastgs-\allowbreak acceleration-\allowbreak method.md} & Nền tảng toán 3DGS + 3 đòn bẩy FastGS + đo FastGS-lite vs 3DGS \\
\texttt{DOCS/colab-t4-guide.md} & Playbook Colab T4, chống OOM, lộ trình CUDA \\
\texttt{DOCS/BOOK/} & Sách hợp nhất 15 chương, từ toán nền tảng tới cost model \\
\texttt{DOCS/BOOK/*\_merge\_figures/} & Biểu đồ minh hoạ công thức các cơ chế Faster-GS mới (fused Adam, 3D filter, Morton, MCMC, random-init) \\
GitHub & \texttt{KietAnhCS/fastgs-lite} \\
\bottomrule
\end{tabular}
\end{table}
\end{frame}

% --- Slide 67 ---
\begin{frame}[c]{}
\centering
{\Huge \textbf{Cảm ơn đã theo dõi!}}

\vspace{6mm}
{\LARGE Hỏi đáp (Q\&A)}

\vspace{10mm}
{\small Digital Twin GS --- FastGS-lite \\ \texttt{github.com/KietAnhCS/fastgs-lite}}
\end{frame}
